% Preâmbulo comum dos gráficos que explicam os métodos (30/10/2026).
% Dados e números: dados/resultados/metodos_ilustrados_2026-10-30/, gerados por
% ferramentas/MetodosIlustrados.mq5 com as funções do próprio robô. Nenhum
% número é digitado aqui: os do texto vêm de valores.tex (\vX para o texto,
% \pX para coordenada).
\usepackage[T1]{fontenc}
\usepackage{amsmath}
\usepackage{mathptmx}  % Times, como as figuras do TG
\usepackage{pgfplots}
\usepgfplotslibrary{groupplots,dateplot,colormaps,fillbetween}
\usetikzlibrary{arrows.meta,positioning}
\pgfplotsset{compat=1.18}
\definecolor{cmod}{RGB}{0,90,181}     % modelo / ajuste
\definecolor{cdes}{RGB}{220,50,32}    % destaque, limite, valor crítico
\definecolor{cdad}{RGB}{120,120,120}  % dado observado
\newcommand{\arq}[1]{../../dados/resultados/metodos_ilustrados_2026-10-30/#1}
\input{\arq{valores.tex}}
% Figuras 08 a 14 (o robô operando): traço por barra da refação de 05/10/2026,
% resumido por ferramentas/scripts/robo_em_acao.sh. Macros \vR e \pR.
\newcommand{\arqr}[1]{../../dados/resultados/robo_em_acao_2026-10-30/#1}
\input{\arqr{valores.tex}}
\definecolor{cari}{RGB}{230,140,0}   % ARIMA puro
\definecolor{cok}{RGB}{0,140,70}     % compra liberada
\usetikzlibrary{shapes.misc}
\newcommand{\passo}[1]{\tikz[baseline=(n.base)]\node[circle,fill=black!75,text=white,inner sep=1pt,font=\footnotesize\bfseries](n){#1};}
% eixo de datas comum às figuras do traço (jul/2019 a jun/2021)
\pgfplotsset{traco/.style={date coordinates in=x, xmin=2019-07-01, xmax=2021-06-30,
    xtick={2019-07-01,2020-01-01,2020-07-01,2021-01-01,2021-06-30},
    xticklabels={jul/2019,jan/2020,jul/2020,jan/2021,jun/2021}}}
\pgfplotsset{
  base/.style={table/col sep=comma, grid=major, grid style={gray!20},
    tick label style={font=\small}, label style={font=\small},
    title style={font=\small\bfseries, align=center},
    every axis plot/.append style={line width=0.9pt},
    /pgf/number format/use comma, /pgf/number format/1000 sep={.}},
}
% faixa vertical nos pregões do traço em que a condição #1 vale, na cor #2,
% do piso #3 ao teto #4 do eixo
\newcommand{\faixa}[4]{%
  \addplot[name path=topo, draw=none, const plot, forget plot] table[x=data, y expr={(#1) ? #4 : #3}]{\arqr{dia_a_dia.csv}};
  \addplot[name path=piso, draw=none, forget plot] coordinates {(2019-07-01,#3) (2021-06-30,#3)};
  \addplot[#2, forget plot] fill between[of=topo and piso];}
